% =====================================================================
%  Chapter 11 — Principle: Fidelity (Respect the Machine)
%  Map: research/chapter-map.md (Chapter 11)
%  Shape: idea → origin → cost → what it opens up. Ideas, not how-to.
%  Rule: the author is never quoted in body text; short quotes only in
%  footnotes. Full statements live in interviews/.
% =====================================================================
\chaptersubtitle{Respect the Machine}
\chapter{Fidelity}
\label{ch:machines}

The C language was designed in the early 1970s for a minicomputer called the
PDP-11, and in a real sense almost every mainstream language since has
inherited its picture of a computer: one processor, one flat expanse of memory,
instructions carried out one after another. The chips on sale today look
nothing like that. They have many cores, units that work on many values at
once, hardware that tags memory and signs pointers to catch attacks, and memory
of different kinds with different properties. Silica is designed to be faithful
to the machines we actually have, not to a portable memory of the ones we
used to have.\footnote{\src{silica/docs/design-principles.md}, line 67.}

\section{The machine as it is}

Fidelity means that what the hardware offers, the language can reach. Different
kinds of memory are distinct in the language because they are distinct in the
chip. Hardware features for safety and speed are treated as part of the world a
language should address, not as exotic extensions bolted on later. Even small
rules are taken from real hardware rather than from tradition.

This is what catches the attention of the bare-metal programmers who do look at
Silica. Every board and every chip exposes behaviours that, in C, they could
reach only with great difficulty, if at all. In Silica they can state the
reality of the device directly: that this memory is of one kind and that memory
of another, that the chip itself enforces atomic access here, that pointers are
audited by the hardware there. Being able to say what the chip actually
guarantees, in the language rather than in comments and conventions, is what
resonates most.\footnote{Interview, 2026-11-13.}

Silica's first home was ARM, but not as a matter of doctrine. ARM chips were
simply what Barney had at work and at home. When the time came for other chips,
a used x86-64 machine was bought and an old Raspberry Pi dug out of storage.
The compiler had been designed so that a new kind of chip needs only a new final
stage, the part that turns a program into that chip's instructions. Ask for a
feature on a chip that lacks it, and the build stops with an error rather than
quietly imitating the feature.

That is fidelity's answer to portability. Portability, the long habit of writing
for the lowest common ground so that a program runs everywhere, matters far less
when specific solutions are easy to produce (Chapter~\ref{ch:transport}). With
AI models to help, developers can do for their own applications what Silica does
for itself, and build each one for the device it runs on. Barney does not regard
the result as bloat. Code built this way is focused on the hardware beneath it
and uses all of it, instead of settling for a bar low enough to reach every
machine.\footnote{Interview, 2026-11-14.}

% TODO-INTERVIEW: A specific board and behaviour someone could finally express
%   in Silica, and their reaction.

\section{Owning the last step}

Most new languages hand their final step---turning a program into machine
instructions---to a shared toolkit, usually LLVM. Silica used LLVM only to get
started, and dropped it as soon as the compiler written in Silica could produce
machine instructions itself. The reason was not speed of compilation. LLVM, like
the languages it serves, thinks in the PDP-11's terms, and the behaviours of
modern chips cannot be expressed efficiently or well through that picture. To
Barney it made no sense to take a broad, expressive language, squeeze it through
a PDP-11-shaped bottleneck, and then hope that optimizers would widen the
running code back out into what the language had said in the first place.

Barney has drawn the comparison on whiteboards many times, and it makes sense to
those who see it (Figure~\ref{fig:hourglass}). The conventional path is an
hourglass. Silica's is a straight channel from the language to the hardware,
with nothing narrower in between.

\begin{figure}[htbp]
  \centering
  % The hourglass and the channel (Chapter 11), after the author's whiteboard
% drawing. Left: an expressive language squeezed through a PDP-11 model, then
% widened again by optimizers. Right: Silica straight down to the hardware.
\begin{tikzpicture}[
    lbl/.style={font=\small, text=taoink},
    note/.style={font=\footnotesize\itshape, text=taoink!75, align=left},
    pathlabel/.style={font=\small\scshape, text=taoink!85},
    arrow/.style={-{Stealth[length=2mm]}, taoink!60, thick},
  ]
  % ---- Left: the hourglass ----
  \begin{scope}
    \fill[taopaper] (-18mm,50mm) -- (18mm,50mm) -- (2.5mm,25mm)
                    -- (18mm,0mm) -- (-18mm,0mm) -- (-2.5mm,25mm) -- cycle;
    \draw[taoink!60, thick] (-18mm,50mm) -- (18mm,50mm) -- (2.5mm,25mm)
                    -- (18mm,0mm) -- (-18mm,0mm) -- (-2.5mm,25mm) -- cycle;
    \node[lbl] at (0,45mm) {Language};
    \node[note, anchor=west] at (7mm,25mm) {PDP-11\\model};
    \node[note] at (0,11mm) {optimizers};
    \node[lbl] at (0,4mm) {Chip};
    \draw[arrow] (0,40.5mm) -- (0,15mm);
    \node[pathlabel] at (0,-6mm) {Conventional};
  \end{scope}
  % ---- Right: the channel ----
  \begin{scope}[xshift=52mm]
    \fill[taopaper] (-13mm,0mm) rectangle (13mm,50mm);
    \draw[taoaccent, very thick] (-13mm,0mm) -- (-13mm,50mm);
    \draw[taoaccent, very thick] (13mm,0mm) -- (13mm,50mm);
    \node[lbl] at (0,45mm) {Silica};
    \node[lbl] at (0,4mm) {Hardware};
    \draw[arrow] (0,40.5mm) -- (0,8.5mm);
    \node[pathlabel] at (0,-6mm) {Silica};
  \end{scope}
\end{tikzpicture}

  \caption{Two paths from language to chip, after Barney's whiteboard drawing.
    The conventional path narrows to a PDP-11 model and relies on optimizers to
    widen it again; Silica's runs straight down to the hardware.}
  \label{fig:hourglass}
\end{figure}

\section{Bare metal first}

Silica describes itself as a language for building systems directly on the
hardware, which also happens to run under ordinary operating systems. Its
history bears that out. Bringing it to a small microcontroller, the project
chose at each turn the harder and more faithful path: generating the chip's own
instructions rather than translating another chip's, running with no vendor
software or operating system underneath, and testing on real boards rather
than simulations. A thousand actors now run on a board with half a megabyte of
memory. More boards followed.

The easier options were turned down for the same reason LLVM was. The PDP-11's
legacy is alive and well in the libraries and tools used to build software for
embedded systems, and taking them would have meant taking the bottleneck with
them. Quicker would have been easier, but better is better, and AI models have
made doing better far more feasible at scale.\footnote{``Quicker is easier. But
better is better.'' Interview, 2026-11-14.}

\section{Honest where control is shared}

On a laptop or a server, the operating system holds some of the control a
systems language would want, deciding how memory behaves and where work runs.
Silica says so, rather than pretending its promises hold equally everywhere
(Chapter~\ref{ch:hidden}). Each kind of machine refuses, when the program is
built, whatever it cannot honestly do, and accepts or refuses programs for the
same reasons on every machine.

\section{Measure first}

Silica's design documents share a habit: describe the present as measured
before proposing a change, and give numbers only where they have been checked.
Several designs in the project's history were withdrawn the moment the
measurements came in.

The habit is older than Silica. Barney once worked in the simulation industry,
on software that had to run very fast, where even a one per cent loss of speed
was wildly unacceptable. That training joined a predilection Barney has had
since youth, for correctness and for beauty.

\section{What it costs}

Faithfulness to real hardware means more work for every new kind of machine,
and fewer shortcuts borrowed from shared toolkits and vendor software. It also
means being honest that a laptop cannot offer everything a bare board can.

\section{What it opens up}

\begin{itemize}
  \item \textbf{Anyone with a need can build.} Asked who should build on
        these boards, Barney's answer is everybody and anybody. The principles
        of solving problems are the same everywhere and in every age, and with
        AI models used well they are enough to build what makes one's own life,
        or others' lives, better. Nothing will be perfect. Failures will happen,
        and the answer is to adjust and move on (Chapter~\ref{ch:failsoft}).
  \item \textbf{One language from the badge to the data centre.} The same
        ideas, and the same programs, can run on a microcontroller and on a
        server.
  \item \textbf{Hardware safety, put to use.} Chip features that catch memory
        attacks can become guarantees of the language rather than optional
        extras few programs use.
  \item \textbf{Freedom from vendor lock-in.} Running directly on the
        hardware, without a vendor's software stack, keeps a device's
        behaviour in its maker's hands.
  \item \textbf{Teaching on real machines.} Inexpensive boards become places
        to learn how computers really work, using the same language
        professionals use.
\end{itemize}

% TODO-INTERVIEW: The simulation industry years: what kind of simulation, and
%   what did it teach beyond measurement?
